%% un-sdg.tex: UN Sustainable Development Goals (FEUP regulations)
%% -------------------------------------------------------

\chapter*{UN Sustainable Development Goals} \label{chap:UnitedNations}
\hspace{\parindent}
The United Nations Sustainable Development Goals (SDGs)~\cite{un_sdgs} provide a global framework to
achieve a better and more sustainable future for all. The framework comprises 17 goals addressing
global challenges such as poverty, inequality, climate change, environmental degradation, peace,
and justice.

This dissertation relates to the SDGs through its contribution to methods aimed at making
centralized multi-robot fleet coordination more robust and efficient. The work can support more
reliable and scalable robotic coordination in industrial and logistics environments through
prioritization heuristics and
parallel planning. The following table identifies the most relevant SDGs, the associated targets,
the contribution of the dissertation, and the corresponding performance indicators.

\begin{description}
\item [SDG 8]
Promote sustained, inclusive and sustainable economic growth, full and productive employment
and decent work for all.

\item [SDG 9]
Build resilient infrastructure, promote inclusive and sustainable industrialization and foster
innovation.

\item [SDG 12]
Ensure sustainable consumption and production patterns.
\end{description}

\begin{center}
\small
\begin{tabular}{|l|l|p{58mm}|p{52mm}|}
\hline
\textbf{SDG} & \textbf{Target} & \textbf{Contribution} & \textbf{Indicators and Metrics} \\
\hline
\hline

\multirow{1}{*}{8}
& 8.2
& The proposed methods can contribute to technological upgrading and operational productivity by improving the robustness and effectiveness of centralized multi-robot fleet coordination in industrial and logistics scenarios.
& Planning success rate; time to first valid solution; solution cost; planning performance in large-scale scenarios. \\
\hline

\multirow{2}{*}{9}
& 9.4
& The proposed planning approach may support more efficient operation of robotic infrastructure by improving the ability to obtain valid coordination plans within a limited planning time.
& Time to first valid solution; number of valid runs; solution cost; scalability across large-scale scenarios. \\
\cline{2-4}

& 9.5
& The dissertation contributes to research and innovation in industrial robotic systems through the design, implementation, and evaluation of prioritization heuristics and a parallel planning framework.
& Number of heuristic families evaluated; comparison with random priority selection; performance across controlled and large-scale scenarios; effect of worker thread configurations. \\
\hline

\multirow{1}{*}{12}
& 12.2
& The proposed improvements indirectly support more efficient use of computational and operational resources by exploring multiple promising priority orderings within the available planning budget.
& Planning success rate; time to first valid solution; solution cost; time to best solution; time to complete the heuristic pool. \\
\hline

\end{tabular}
\end{center}